% Kamp III.2, Lemma 2: limiting finite patterns over all linear orders.

\subsubsection{Limiting past and future patterns}
\label{sec:p09-kamp-iii2-lemma2}

The next lemma expresses finite patterns that can be realized inside every
past or future interval adjoining the current point.  Kamp's Lemma~2 states
expressibility of the limiting classes over \(\mathsf{Lin}\); its proof
occupies printed pp.~70--73 \cite[ch.~III, lemma~2, pp.~66--73]{Kamp1968}.
The construction below supplies both directions and every possible position
at which a provisional finite chain reaches the current point.

\begin{hbdefinition}{def:p09-kamp-limit-operators}{The two strict limit operations}
For a temporal formula \(\eta\), use the following abbreviations in
\(\mathsf{TL}_1\):
\[
\begin{aligned}
 H'(\eta)&\defeq S(\neg(\eta\land\neg\eta),\eta),&
 G'(\eta)&\defeq U(\neg(\eta\land\neg\eta),\eta),\\
 P'(\eta)&\defeq\neg H'(\neg\eta),&
 F'(\eta)&\defeq\neg G'(\neg\eta).
\end{aligned}
\]
These are Kamp's derived operations, with the strict witness and open-interval
argument order already fixed for \(S,U\)
\cite[ch.~III, defs.~10--11, pp.~57--59]{Kamp1968}.
For a first-order formula \(\theta\) whose sole free variable is \(w\),
write
\[
\begin{aligned}
 \partial^-\theta(w)&\defeq
   \forall x\bigl(x<w\to\exists z\,
        (x<z\land z<w\land\theta[z/w])\bigr),\\
 \partial^+\theta(w)&\defeq
   \forall x\bigl(w<x\to\exists z\,
        (w<z\land z<x\land\theta[z/w])\bigr).
\end{aligned}
\]
The bound variables \(x,z\) are fresh.  The symbols \(\partial^-,\partial^+\)
name these formula operations; they do not denote analytic derivatives.
The annex proves that if \(\eta\) expresses \(\theta(w)\), then
\(P'(\eta)\) and \(F'(\eta)\) express these two formulas.

The universal conditions include empty past and future domains.  In
particular, \(\partial^-\theta\) holds at a first point and
\(\partial^+\theta\) holds at a last point, whatever \(\theta\) says.
They impose no condition on \(\theta\) at the current point itself.
\end{hbdefinition}

\begin{hbdefinition}{def:p09-kamp-limiting-classes}{Limiting basic-pattern classes}
Keep the basic classes and occurrence words of
\cref{def:p09-kamp-basic-classes,def:p09-kamp-basic-word}.
For each unary predicate letter \(A\), the base formulas are
\(\partial^-A(w)\) and \(\partial^+A(w)\).  Set
\[
 \operatorname{Lp}_0(w)\defeq\{\partial^-A(w):A\text{ unary}\},
 \qquad
 \operatorname{Lf}_0(w)\defeq\{\partial^+A(w):A\text{ unary}\}.
\]
For every \(n\geq1\), define
\[
\begin{split}
 \operatorname{Lp}_n(w)\defeq\bigl\{
  &\forall x\bigl(x<w\to\exists u\,
     (x<u\land u<w\land\exists v\,
       (u<v\land v<w\land\phi))\bigr):\\
  &\hspace{35mm}\phi\in\operatorname{Bp}_n(u,v)\bigr\},\\
 \operatorname{Lf}_n(w)\defeq\bigl\{
  &\forall x\bigl(w<x\to\exists u\,
     (w<u\land u<x\land\exists v\,
       (u<v\land v<x\land\phi))\bigr):\\
  &\hspace{35mm}\phi\in\operatorname{Bp}_n(u,v)\bigr\}.
\end{split}
\]
All displayed endpoint names are distinct and all other bound variables are
fresh.  Thus each resulting formula has exactly \(w\) free.  The positive
index counts the \(n\) interval pieces of its constituent basic pattern;
the base index zero concerns one predicate letter.  There is no retention
of earlier stages inside a later stage.  Put
\[
 \operatorname{Lp}(w)\defeq\bigcup_{n\geq0}\operatorname{Lp}_n(w),
 \qquad
 \operatorname{Lf}(w)\defeq\bigcup_{n\geq0}\operatorname{Lf}_n(w),
\]
and let \(\operatorname{Lp},\operatorname{Lf}\) be their unions over the
choice of free variable.  Both source base definitions test their predicate
at the quantified witness point.  Moving that test to the current point
changes the condition, as the handbook comparison in
\cref{def:p09-kamp-current-point-comparison} demonstrates.
The higher-stage indexing is stated by interval count throughout
\cite[ch.~III, pp.~63, 66--67, 70--72]{Kamp1968}.
\end{hbdefinition}

\begin{hbdefinition}{def:p09-kamp-limit-translation}{The finite temporal construction}
Let
\(\sigma=(A_0,B_1,A_1,\ldots,B_n,A_n)\) be an occurrence word.
For \(n\geq1\), let \(E^-_\sigma\) be the formula \(L_n\) and
\(E^+_\sigma\) the formula \(R_0\) constructed in
\cref{thm:p09-kamp-iii2-lemma1}.  They say respectively that a realization
of the whole word ends or begins at the current point.  For \(n=0\),
set \(E^-_{(A_0)}=E^+_{(A_0)}=\widehat A_0\).

For \(n\geq1\), the suffix and prefix words are
\[
 \sigma^-=(A_1,B_2,A_2,\ldots,B_n,A_n),\qquad
 \sigma^+=(A_0,B_1,A_1,\ldots,B_{n-1},A_{n-1}).
\]
They each have one fewer interval; at \(n=1\) they are the one-letter
words \((A_1)\) and \((A_0)\).  Define recursively
\[
\begin{aligned}
 \Lambda^-_{(A)}&\defeq P'(\widehat A),&
 \Lambda^+_{(A)}&\defeq F'(\widehat A),\\
 \Lambda^-_\sigma&\defeq
            P'(E^+_\sigma)\land\Lambda^-_{\sigma^-},&
 \Lambda^+_\sigma&\defeq
            F'(E^-_\sigma)\land\Lambda^+_{\sigma^+}
                    \qquad(n\geq1).
\end{aligned}
\]
This is a finite construction: each recursive call removes exactly two
predicate occurrences and one interval.  All its leaves are proposition
letters already occurring in the input word.  The abbreviations \(H',G',P',F'\)
expand into negation, conjunction and the original \(S,U\).
\end{hbdefinition}

\begin{hblemma}{thm:p09-kamp-iii2-lemma2}{proved=annex:p9.kamp.iii2.lemma2}{Kamp's limiting past and future lemma}
Every member of \(\operatorname{Lp}\cup\operatorname{Lf}\) is expressible
over \(\mathsf{Lin}\).  More explicitly, a member of
\(\operatorname{Lp}_n(w)\), with constituent occurrence word \(\sigma\),
is expressed by \(\Lambda^-_\sigma\); a member of
\(\operatorname{Lf}_n(w)\) is expressed by \(\Lambda^+_\sigma\).
At stage zero, the constituent word is its single predicate letter.

These assertions use the named-variable \(D\)-translation of
\cref{def:p09-kamp-basic-expressibility}, with capture-avoiding replacement
of \(t_0\) by \(w\).  They hold on arbitrary nonempty strict linear
orders, with arbitrary interpretations of the unary predicates.  Density,
endpointlessness and order completeness are not hypotheses
\cite[ch.~III, lemma~2, pp.~70--73]{Kamp1968}.
\end{hblemma}

\begin{hbexample}{def:p09-kamp-limit-one-piece}{One interval and the required second conjunct}
For the word \((A,B,C)\), the two formulas are
\[
\begin{aligned}
 \Lambda^-_{(A,B,C)}
   &=P'(\widehat A\land U(\widehat C,\widehat B))
                                      \land P'(\widehat C),\\
 \Lambda^+_{(A,B,C)}
   &=F'(\widehat C\land S(\widehat A,\widehat B))
                                      \land F'(\widehat A).
\end{aligned}
\]
The first conjunct alone permits a provisional chain whose other endpoint
lies beyond the current point.  For example, on \(\mathbb R\) at zero,
take \(A=(-\infty,0)\), \(B=\mathbb R\), and \(C=\{1\}\).
Then \(P'(\widehat A\land U(\widehat C,\widehat B))\) is true, while
the limiting past pattern and \(P'(\widehat C)\) are false.
The future counterpart takes \(A=\{-1\}\), \(B=\mathbb R\), and
\(C=(0,\infty)\): its first future conjunct is true and its second
conjunct and limiting future pattern are false.
\end{hbexample}

\begin{hbexample}{def:p09-kamp-limit-endpoints}{Endpoints, adjacency and an extremal crossing}
At a first point every limiting past formula above is true; at a last point
every limiting future formula is true.  If a point has an immediate
predecessor, every limiting past formula is false.  If it has an immediate
successor, every limiting future formula is false.  In particular both
families hold at the point of a singleton order, independently of the
predicate interpretation.

The last possible crossing must be included in the past splice.  With all
predicates true on \(\mathbb R\), the two-interval chain
\((-3,-1,0)\) begins above the cut \(-4\) but reaches the current point
zero.  A suffix realization \((-1/2,-1/4)\) lies above its last strictly
past point \(-1\).  Splicing gives \((-3,-1,-1/4)\), wholly below zero.
The crossing is between the last two original points and includes equality
of the old last point with the current point.  The proof covers this case
and its future counterpart explicitly.
\end{hbexample}

\begin{hbexample}{def:p09-kamp-current-point-comparison}{A current-point condition and a strict limit}
Replacing the witness-point predicate \(A(z)\) by \(A(w)\) produces a
different formula.  For this additional handbook comparison, write
\[
 \operatorname{CurF}_A(w)\defeq
  \forall x\bigl(w<x\to\exists z\,
                  (w<z\land z<x\land A(w))\bigr).
\]
This also is expressible over \(\mathsf{Lin}\).  Put
\(\tau_A\defeq\neg(\widehat A\land\neg\widehat A)\); its expression is
\[
 \neg U(\tau_A,\tau_A)\ \lor\
                  (\widehat A\land F'(\tau_A)),
\]
where disjunction abbreviates \(\neg(\neg\alpha\land\neg\beta)\).
On \(\mathbb R\) at zero with \(A=\{0\}\), this current-point condition is
true and the source's witness-point base \(\partial^+A\) is false.
Thus a condition on \(A\) at the current point cannot replace the
witness-point clause in the limiting definition.  The annex proves this
comparison as well as the source-based limiting construction.
\end{hbexample}

\paragraph{The next dependency.}
The next source result is Lemma~3 on printed p.~73.  It concerns the closure
of current-point, basic and limiting classes under the Boolean operations
and proper predicate substitution, called \(E\) in the source.  Its
substitution step and the subsequent interval lemmas remain separate work.
Lemma~3 still uses \(\mathsf{Lin}\); only after it does Kamp change the
default convention to \(\mathsf{Com}\).  Lemma~2 alone does not establish
the full expressibility result of Theorem~III.2.
